\documentclass[12pt,a4paper]{article}
\usepackage[utf8]{inputenc}
\usepackage[italian]{babel}
\usepackage{amsmath, amssymb}
\usepackage{graphicx}
\usepackage{geometry}
\geometry{a4paper, top=3cm, bottom=3cm, left=2.5cm, right=2.5cm}
\usepackage{hyperref}
\usepackage{listings}
\usepackage{xcolor}

\title{\textbf{Relazione di Progetto: Controllo Autonomo di uno Sciame di Droni Sottomarini tramite FMI e Reti Neurali Evolutive}}
\author{Massimiliano}
\date{\today}

\begin{document}

\maketitle

\tableofcontents
\newpage

\section{Introduzione}
Il presente progetto ha l'obiettivo di sviluppare, simulare e validare un sistema di controllo autonomo per uno sciame di droni sottomarini. L'architettura del sistema si basa sull'integrazione di modelli fisici continui, sviluppati e validati in \textit{OpenModelica}, con un'intelligenza artificiale addestrata tramite \textit{Evolution Strategies} in un ambiente C custom. Infine, è stata implementata una logica di missione ad alto livello tramite una Macchina a Stati Finiti (FSM) che gestisce autonomamente l'esplorazione spaziale e i cicli di ricarica energetica, garantendo la sopravvivenza a lungo termine della flotta.

\section{Modellazione Fisica in OpenModelica}
La base fisica della simulazione è stata sviluppata interamente nell'ambiente OpenModelica. È stato creato un modello che descrive la cinematica, la dinamica e il consumo energetico di un drone sottomarino.

\subsection{Dinamica e Cinematica}
Il drone è stato schematizzato come un corpo rigido soggetto alle forze di propulsione e alla resistenza idrodinamica dell'acqua. Le variabili di stato principali del sistema includono la posizione spaziale $(x, y, z)$ e la velocità vettoriale $(v_x, v_y, v_z)$. 
Le equazioni del moto implementate nel modello Modelica tengono conto di un coefficiente di attrito viscoso (drag) proporzionale alla velocità, e della spinta generata dai propulsori $(u_x, u_y, u_z)$, che funge da variabile di input per il controllore. Sono stati inoltre imposti dei limiti fisici di spinta per simulare le limitazioni dei rotori reali.

\subsection{Modellazione Energetica}
Per rendere la simulazione più aderente alla realtà, è stato modellato un sistema di gestione della batteria. Lo stato di carica (SOC) decresce in funzione del tempo e dello sforzo dei propulsori. In particolare, è stata implementata una logica secondo cui la batteria si ricarica automaticamente quando il drone si trova a riposo (spinta pari a zero) e in uno specifico stato di "recharging".

\subsection{Esportazione tramite FMI (Functional Mock-up Interface)}
Il modello OpenModelica è stato successivamente esportato come \textit{Functional Mock-up Unit} (FMU) aderente allo standard FMI 2.0 per la \textit{Co-Simulation}. Questo ha permesso di separare la simulazione fisica (che rimane incapsulata e risolta internamente dall'FMU con il proprio solver) dalla logica di controllo e addestramento scritta in C, permettendo prestazioni di esecuzione eccezionali necessarie per il machine learning.

\section{L'Ambiente di Simulazione in C}
Per addestrare e testare i droni, è stato scritto un ambiente di simulazione da zero in linguaggio C, il quale si interfaccia con le istanze FMU tramite la libreria \textit{FMILibrary}.

\subsection{Generazione Procedurale della Mappa}
La mappa sottomarina è generata casualmente a ogni episodio per garantire che la Rete Neurale non "impari a memoria" un percorso specifico, ma generalizzi l'abilità di navigazione.
L'ambiente genera un fondale (mare) su $Y=0$, e dei limiti di perimetro sui lati e sul soffitto. Inoltre, vengono istanziati molteplici ostacoli sferici (o semisfere ancorate al fondale) con raggio e posizioni pseudo-casuali, garantendo tuttavia che non si sovrappongano al punto di spawn del drone.

\subsection{Sensori LIDAR (Raycasting)}
Per permettere alla Rete Neurale di "vedere" gli ostacoli, è stato implementato un sistema LIDAR omnidirezionale. Tramite l'algoritmo della \textit{Sfera di Fibonacci}, vengono generati $N = 64$ raggi distribuiti uniformemente nello spazio. Per ciascun raggio, il motore C calcola l'intersezione analitica con tutte le geometrie della mappa (limiti fisici, ostacoli sferici, e persino altri droni), restituendo la distanza dell'ostacolo più vicino.

\section{Intelligenza Artificiale: Reti Neurali ed Evolution Strategy}
Il "cervello" di ogni drone è una Rete Neurale \textit{Feed-Forward} (MLP).

\subsection{Architettura della Rete Neurale}
L'architettura adottata prevede:
\begin{itemize}
    \item \textbf{Input Layer (68 neuroni):} 64 valori derivanti dalle distanze normalizzate del Lidar, 3 valori che rappresentano il vettore direzionale normalizzato verso il target $(dx, dy, dz)$, e 1 valore che rappresenta lo stato della batteria.
    \item \textbf{Hidden Layer (32 neuroni):} Utilizza una funzione di attivazione \textit{tanh} per catturare relazioni non lineari.
    \item \textbf{Output Layer (3 neuroni):} Utilizza l'attivazione \textit{tanh} per fornire l'input di spinta $[-1, 1]$ per i tre assi, successivamente scalato per la forza massima del propulsore.
\end{itemize}

\subsection{Addestramento (OpenAI Evolution Strategy)}
Considerata la difficoltà di calcolare gradienti in un ambiente di simulazione \textit{black-box} (FMU), è stato implementato l'algoritmo \textit{Evolution Strategy} (basato sul paper di OpenAI).
Il processo è iterativo (oltre 1000 generazioni):
\begin{enumerate}
    \item Viene generata una "popolazione" di variazioni pseudo-casuali (jittering) attorno ai pesi attuali della rete neurale.
    \item Ogni variante viene testata simulando un intero episodio (eseguendo step sull'FMU).
    \item La \textit{Fitness} viene valutata premiando i droni che si avvicinano al target e penalizzando pesantemente collisioni con ostacoli o l'esaurimento della batteria.
    \item I pesi della rete vengono aggiornati facendo una media pesata delle variazioni proporzionalmente al punteggio di fitness ottenuto.
\end{enumerate}
Il modello con le prestazioni migliori viene salvato in un file binario (\texttt{best\_model.bin}).

\section{Il Sistema di Missione Dinamica}
Una volta addestrata la Rete Neurale a evitare ostacoli e raggiungere target in ambienti sconosciuti, è stato costruito un livello di controllo logico superiore, il \textit{Mission System}, programmato interamente in C.
Questa architettura permette allo sciame di operare su scenari continui di durata indefinita.

\subsection{Macchina a Stati Finiti (FSM)}
Ogni drone è governato da una FSM con tre stati:
\begin{enumerate}
    \item \textbf{STATE\_SEEK\_TARGET:} Il drone esplora l'ambiente. Il sistema gli fornisce una coordinata tridimensionale casuale all'interno di un raggio $I\_RADIUS = 20.0$ metri, validata per essere lontana da ostacoli. Quando il drone raggiunge la coordinata, gli viene assegnato un nuovo target, simulando un pattugliamento continuo. L'input per il moto in questa fase è interamente affidato alla Rete Neurale pre-addestrata.
    \item \textbf{STATE\_RETURN\_BASE:} Quando il sensore della batteria scende sotto la soglia critica del $30\%$, la FSM interviene sostituendo il target di esplorazione con le coordinate originarie di spawn (la Base di Ricarica). Il drone si affida sempre alla Rete Neurale per compiere il tragitto a ritroso, evitando gli ostacoli lungo la via del ritorno.
    \item \textbf{STATE\_RECHARGING:} Una volta arrivato alla base (distanza $< 5.0m$), i propulsori vengono disattivati (spinta = $0.0$). L'FMU simula fisicamente la ricarica. Il drone attende in questo stato finché la batteria non raggiunge il $90\%$, al termine del quale la FSM ricomincia il ciclo \textit{SEEK\_TARGET}.
\end{enumerate}

\section{Visualizzazione e Strumentazione 3D}
Per validare visivamente il comportamento fisico, le collisioni e la logica di missione, sono stati sviluppati script in Python basati su \textit{Matplotlib}.
Il motore C, sia durante l'addestramento sia in fase di missione, esporta lo stato dell'ambiente (posizioni degli ostacoli, bounds) e la telemetria continua ad alta frequenza in formati \texttt{.csv}.
Il visualizzatore 3D (\texttt{visualize\_mission.py}) legge questi output per generare un'animazione interattiva fluida che mostra in tempo reale:
\begin{itemize}
    \item La struttura della mappa (mare, muri, ostacoli).
    \item Le traiettorie e i punti correnti dei droni.
    \item Lo spostamento dinamico delle "stelle" obiettivo a mano a mano che vengono scoperte.
    \item Il riepilogo a schermo degli stati della batteria e della FSM (Seek, Return, Charge).
    \item I raggi Lidar emessi, evidenziando le collisioni imminenti.
\end{itemize}

\section{Conclusioni}
Il progetto dimostra con successo l'integrazione di paradigmi avanzati dell'ingegneria del software e dell'IA. Il disaccoppiamento tra il modello fisico matematico puro (OpenModelica/FMI) e la logica decisionale e di training (C/Neural Networks) garantisce una straordinaria flessibilità. L'addestramento evolutivo ha consentito l'emergenza di un comportamento intelligente generalizzato di navigazione \textit{obstacle-avoidance}, solido abbastanza da essere incapsulato e gestito da una logica a stati per missioni persistenti (energy management) nel mondo reale.

\end{document}
